\exercisegroup

\begin{enumerate}
\def\labelenumi{\arabic{enumi})}
\item
  Let X and Y be topological spaces, and let \(f: X \to Y\) be a continuous map. Let \((T, p)\) be a Y-space. The canonical descent datum \(\tau\) on the X-space \(Z = X \times_{Y} T\) is effective (IV, p. 385, prop. 1), the canonical map \(Z/R_{\tau} \to T\) is injective and continuous. It is not, however, necessarily surjective, nor strict.
\item
  Let Y be a topological space obtained by gluing spaces \(X_{i}\). Let X be the topological space that is the disjoint union of the family \((\mathrm{X}_{i})\) and let \(f: X \to Y\) be the canonical map (TG, I, p. 16). For every i, set \(Y_{i} = f(\mathrm{X}_{i})\) and let \(f_{i}: X_{i} \to Y_{i}\) be the map deduced from f by passing to subspaces. The map \(f_{i}\) is bijective and continuous. If the map \(f_{i}\) is not a homeomorphism, the canonical descent datum on X relative to f is not effective. For an example, cf. TG, I, p. 94, exerc. 15.
\item
  Let B denote the set of points \((x,y)\) of the square \([0,1]^{2}\) satisfying x=0, or x=1, or y=0, or y=1, or there exists \(n\in N\) such that nx=1. Let \(B_{1}\) and \(B_{2}\) the subsets of B consisting of the points \((x,y)\) of B for which \(y\leqslant\frac{1}{2}\) and \(y\geqslant\frac{1}{2}\) respectively.
\end{enumerate}

Construct a B-space \((\mathrm{E}, p)\) which is not a covering but such that the \(B_{1}\) -space \((\mathrm{E}_{\mathrm{B}_{1}}, p_{\mathrm{B}_{1}})\) and the \(B_{2}\) -space \((\mathrm{E}_{\mathrm{B}_{2}}, p_{\mathrm{B}_{2}})\) are coverings.

\begin{enumerate}
\def\labelenumi{\arabic{enumi})}
\setcounter{enumi}{3}
\tightlist
\item
  A topological space is said to be locally simply connected by arcs if every point has a base of neighborhoods that are simply connected by arcs.
\end{enumerate}

\begin{enumerate}
\def\labelenumi{\alph{enumi})}
\item
  Give an example of a topological space that is simply connected by arcs but is not locally simply connected by arcs.
\item
  Let X be a topological space and let G be a discrete group acting properly on X. If X is locally simply connected by arcs, then so is X/G.
\end{enumerate}

\begin{enumerate}
\def\labelenumi{\arabic{enumi})}
\setcounter{enumi}{4}
\item
  Let G be a Lie group acting properly on a loopable and completely regular topological space X. Set Y = X/G and denote \(f: X \to Y\) the canonical map. Prove that the canonical groupoid morphism \(\varpi(X)/G \to \text{Coeg}(f)\) is an isomorphism. (Use LIE, IX, §9, exerc. 17 to adapt the proof of Theorem 3 of IV, p. 403.)
\item
  Let X be a path-connected topological space, let G be a group acting continuously on X, and let x be a point of X such that \(g \cdot x = x\) for all \(g \in G\) . Denote by p the canonical map from X to X/G. Let H be a groupoid and let \(\varphi: \varpi(X) \to H\) a morphism of groupoids. Suppose that for every path \(c\) on \(\mathrm{X}\) and every \(g\in \mathrm{G}\) , we have \(\varphi ([g\cdot c]) = \varphi ([c])\) and that for every loop \(c\in \Omega_x(\mathrm{X})\) , \(\varphi ([c]) = e_{\varphi (x)}\) .
\end{enumerate}

Let \(c_{1}\) and \(c_{2}\) be paths in X such that \(p \circ c_{1}\) and \(p \circ c_{2}\) have the same origin and the same endpoint. Prove that we have \(\varphi([c_{1}]) = \varphi([c_{2}])\) .

\begin{enumerate}
\def\labelenumi{\arabic{enumi})}
\setcounter{enumi}{6}
\tightlist
\item
  Let X be a loopable topological space and let G be a discrete group acting properly on X; we denote by \(m\colon G \times X \to X\) the action of G, \(\mathrm{pr}_2\colon G \times X \to X\) the second projection, and \(f\colon X \to X / G\) the canonical surjection.
\end{enumerate}

\begin{enumerate}
\def\labelenumi{\alph{enumi})}
\item
  Prove that the groupoid morphism \(\varpi(f)\) is surjective.
\item
  Let H be a groupoid and let \(\varphi\colon\varpi(X)\to H\) a morphism of groupoids such that \(\varphi\circ\varpi(m)=\varphi\circ\varpi(\mathrm{pr}_{2})\) .
\end{enumerate}

Let \(c_{1}\) and \(c_{2}\) be paths in X such that the paths \(f \circ c_{1}\) and \(f \circ c_{2}\) in X/G are strictly homotopic. Prove that one has \(\varphi([c_{1}]) = \varphi([c_{2}])\) .

\begin{enumerate}
\def\labelenumi{\alph{enumi})}
\setcounter{enumi}{2}
\tightlist
\item
  Prove that there exists a unique groupoid morphism \(\varphi'\) from \(\varpi(\mathrm{X/G})\) into H such that \(\varphi = \varphi' \circ \varpi(f)\). Deduce from this a new proof of Theorem 3 of IV, p. 403.
\end{enumerate}

